\section{A moving family forces a lower bound of 9/4}\label{sec:motion}

A tensor depending on a scalar parameter can have a different
price-minimizing placement at each parameter value.
Section~\ref{sec:sync} will construct these placements on a disk. We
show that a nonconstant family satisfying the disk comparisons forces
square matrix multiplication to have exponent at least \(9/4\). The first step is
geometric. At the depth where a translation first moves a point, a
rescaled neighborhood contains a tree in which two branches meet at
the valuation of the difference of their parameters. The second step
uses prices on four equal branches of that tree. With the fourth tip
at a generic parameter, a label shared by all four tips has price at
most \(3/2\) times the price of a wire spanning the diameter, to
first order.
Section~\ref{sec:boundary} uses collisions between the tips to pass
from this estimate to a three-site comparison for matrix multiplication rank.

Let \(F\) be algebraically closed. Take an uncountable set \(\mathcal U\)
of indeterminates, and put
\[
 K_0=\operatorname{Frac}\bigl(F(\mathcal U)[t^{\mathbb Q}]\bigr),
 \qquad \mathcal O=\{z:\nu(z)\ge0\},\qquad J=\{z:\nu(z)>0\},
\]
with the lowest-exponent valuation and \(\nu(0)=+\infty\).
An indeterminate is \emph{fresh} for finite data in \(K_0\) if it
does not occur in those data. Consider a family of sites \(x_z\),
\(z\in J\), in a spectral space over \(F\), subject to the following
conditions.
\begin{enumerate}
\item Adding the same \(w\in J\) to all indices, and permuting
  \(\mathcal U\), preserve the prices of all tensors over \(F\).
\item For a nonzero tensor on any finite list of slots, let \(P(c)\)
  be its price at the sites \(x_{a_i+b_i c}\), where \(a_i,b_i\in J\).
  With \(u\) fresh, \(P(c)\le P(u)\) for \(c\in\mathcal O\).
  For \(c_j\in\mathcal O\) of distinct residues and positive rational
  lengths \(s_j,s_\infty\), the same test satisfies
  \begin{equation}\label{eq:motion-star}
   \sum_j\frac{P(c_j+t^{s_j}u)-P(u)}{s_j}
       +\frac{P(t^{-s_\infty}u)-P(u)}{s_\infty}\ \ge\ 0
  \end{equation}
  whenever the indices at the outer end lie in \(J\). Here \(u\) is
  fresh also for the star data. Thus the price at the generic point
  of the disk is at most the weighted mean of its end prices, with
  weights proportional to the reciprocals of the corresponding
  lengths. The two comparisons are the ordinary-price versions of
  specialization and the submean inequality in Lemma~\ref{lem:star}.
\item Some two sites have positive distance, and
  \(d(x_0,x_z)\to0\) as \(\nu(z)\to+\infty\).
\end{enumerate}
The inequalities apply in particular to distances, by testing a
two-site wire and dividing its price by the logarithm of its size.

\begin{theorem}[Rigidity of analytic motion]\label{thm:motion}
 If such a family exists, square matrix multiplication over $F$ has exact rank at least $n^{9/4}$ for each side $n$.
\end{theorem}

\subsection{The threshold for displacement}

To find a scale at which motion persists after rescaling, fix a point
under all translations of sufficiently high valuation. Its orbit
starts to move when the valuation falls below the threshold for that
point. We shall choose the point so that its displacement has a
strictly negative left derivative at the threshold.

By translation invariance \(x_w\ne x_0\) for some \(w\in J\). Choose
rational \(0<b<\nu(w)\) and replace \(x_z\) by \(x_{t^b z}\). This
still gives a nonconstant family satisfying all three conditions.
Specialization of the wire \(d(x_0,x_{t^bc})\) bounds the radius of
the new family by its value at a fresh \(u\). Thus we can work in the
bounded closed convex hull \(X\) of the new sites, taken in a complete
CAT(0) spectral extension. By the geodesic determination of prices in
Section~\ref{sec:prices}, the symmetries act on \(X\) by
price-preserving isometries. Write \(g_w\) for the action of translation.

We first extend the affine comparisons in Condition~2 to sites
\(g_{a_i+b_i c}A_i\) with \(A_i\in X\). If \(A_i\) is obtained by a
finite sequence of geodesic interpolations between original sites,
every segment in its construction undergoes the same translation as
\(A_i\). Its length and the two distances specifying the interpolation
point are constant across the comparison. Lemma~\ref{lem:geodesic}
therefore extends the star inequality through each interpolation. It
extends specialization in the same way, with the specialized placement
as center and the generic placement as its sole end. Now approximate
arbitrary \(A_i\) by such points and apply the transport inequality
\eqref{eq:prices-transport}. The approximating sequences for these
finitely many sites involve only countably many indeterminates, so one
choice of \(u\) serves throughout. If every \(A_i\) is invariant under
permutations, permuting indeterminates while fixing the finite test
data makes the comparison available at any fresh \(u\).

For \(r>0\) rational and \(A\in X\), take \(u\) outside a countable
support for \(A\), such as the variables in a sequence of finite
geodesic constructions converging to \(A\). Specialization gives
\begin{equation}\label{eq:motion-R}
 R(A,r):=\sup_{\nu(w)\ge r}d(A,g_w A)
          =d(A,g_{t^r u}A).
\end{equation}
For any given \(w\) with \(\nu(w)\ge r\), specialization bounds its
displacement by the generic displacement at a variable fresh for both
\(A\) and \(w\). Permuting that variable to \(u\) while fixing a
countable support for \(A\) gives the same bound for every \(w\).
Conversely, \(t^r u\) itself belongs to the subgroup. Apply
\eqref{eq:motion-star} to the distance \(d(A,g_{t^rc}A)\) with one
inner end at zero. For positive rational \(r-s,r,r+b\) with
\(s,b>0\), this gives
\[
 \frac{R(A,r)-R(A,r-s)}{s}
 \ \le\ \frac{R(A,r+b)-R(A,r)}{b}.
\]
Thus \(R(A,\cdot)\) is convex. It lies between zero and
\(\operatorname{diam}X\), so it extends continuously to \(r>0\) with
a Lipschitz bound on each compact subinterval independent of \(A\).
Squeezing the supremum in \eqref{eq:motion-R} between its values at
rational depths identifies it with this extension at real depths.

The common fixed set \(X_r=\{A:R(A,r)=0\}\) of the subgroup
\(\nu(w)\ge r\) is closed and convex. The circumcenter of any orbit of
this subgroup belongs to \(X_r\). Since the orbit of \(A\) contains
\(A\) and fits in the ball of radius \(R(A,r)\) about \(A\), this gives
\(d(A,X_r)\le R(A,r)\). In particular,
\(f(r)=d(x_0,X_r)\) is positive whenever the subgroup moves \(x_0\),
decreases with \(r\), and tends to zero by Condition~3. For
\(0<h<r\), the projection \(A_r\) of \(x_0\) to \(X_r\) satisfies
\[
 0\le f(r-h)-f(r)\le R(A_r,r-h)-R(A_r,r).
\]
The common Lipschitz bound makes \(f\) locally Lipschitz. Its strict
decrease somewhere gives a depth \(r_*\) with \(f'(r_*)<0\). Fix
\(A=A_{r_*}\). Every permutation of the indeterminates fixes \(x_0\)
and preserves \(X_{r_*}\), so it fixes this unique projection.
The last display bounds the left derivative of the convex function
\(R(A,\cdot)\) at \(r_*\) by \(f'(r_*)\). To the right of \(r_*\) it
vanishes identically. Writing \(-k<0\) for the left derivative gives
\begin{equation}\label{eq:motion-kink}
 R(A,r_*+h)=k(-h)_+ +o(|h|).
\end{equation}

\subsection{Rescaling at the threshold}

Take rational \(\varepsilon\downarrow0\) and choose rational
\(r_0=r_*-\varepsilon+o(\varepsilon)\). In the pointed spectral
ultralimit at \(A\), rescale distances and prices by \(2/(k\varepsilon)\).
For each \(z\in K_0\) take the point \(L_z\) represented by
\[
   g_{t^{r_0}z(t^\varepsilon)}A ,
\]
where the substitution fixes \(F(\mathcal U)\). For fixed \(z\) the
translation index is in \(J\) eventually and the rescaled displacement
is bounded by \eqref{eq:motion-kink}. Finite initial terms can be
chosen arbitrarily. Additivity of the index substitution gives
translation invariance for all \(z\in K_0\). Permutation invariance
follows since every permutation fixes \(A\) and commutes with the
substitution. To transfer the star comparison for arbitrary affine
expressions in \(c\), observe that their coefficients become
\(t^{r_0}a_i(t^\varepsilon),t^{r_0}b_i(t^\varepsilon)\) before passing
to the limit. The centers \(c_j(t^\varepsilon)\) still have distinct
residues, each star length is multiplied by \(\varepsilon\), and all
the tested indices lie in \(J\) eventually. Thus the common factors
in \eqref{eq:motion-star} cancel and the inequality passes to the
ultralimit. The same argument applies to specialization.

At a generic parameter the limiting distance follows directly from
\eqref{eq:motion-R} and \eqref{eq:motion-kink},
\[
 d(L_0,L_{t^ru})=2(1-r)_+
 \qquad(r\in\mathbb Q).
\]
The same formula holds at every parameter of that valuation. More
precisely,
\begin{equation}\label{eq:motion-leaves}
 d(L_z,L_{z'})=2(1-\nu(z-z'))_+ .
\end{equation}
For distinct indices write \(z'-z=t^r v\) with \(\nu(v)=0\).
Translation and specialization give the upper bound, which is already
equality for \(r\ge1\). If \(r<1\), test \(d(L_0,L_{t^rc})\) on a star
with two inner ends centered at zero and at \(v\). The slope on the
zero arm of any length \(0<a<1-r\) is \(-2\). On the outer arm it is
\(+2\), by the generic distance formula. These slopes cancel in
\eqref{eq:motion-star}, leaving a nonnegative slope on the \(v\)-arm.
Specialization makes that slope nonpositive, so its end distance
equals the generic value \(2(1-r)\). Choose the length \(b\ge1-r\) on
this arm. The upper bound already proved identifies
\(L_{t^r(v+t^bu)}=L_{t^rv}\), which gives equality in
\eqref{eq:motion-leaves}.

\subsection{Geodesics of the disk tree}

The distance formula places the leaves. To locate the branch points,
we prove that the midpoint of two leaves lies on each geodesic joining
one of them to a third leaf at least as far away as the other.

Let \(B=L_b\), \(C=L_{b+v}\), \(d(B,C)=2h>0\),
and let \(m\) be their midpoint. Take a star with inner ends \(0,1\),
each of length \(h\), and outer length \(s\). On \([B,C]\), prescribe
the point \(m\) at the center and the outer end, and \(B,C\) at the
two inner ends. Its distances to each endpoint are harmonic for the
weights proportional to the inverse end lengths. Lemma~\ref{lem:geodesic}
therefore permits comparing its distance to \(L_{b+cv}\). At the inner
ends this distance is zero by \eqref{eq:motion-leaves}. At the outer
end it is at most \(3h+2s\) by going through \(B\). Consequently the
generic center distance \(J_*=d(m,L_{b+uv})\) satisfies
\[
  0\le-\frac{2J_*}{h}+\frac{3h+2s-J_*}{s}.
\]
Sending \(s\) to infinity yields \(J_*\le h\), and specialization
gives the same bound throughout the unit disk. More generally,
\begin{equation}\label{eq:motion-midpoint}
 d(m,L_{b+cv})\le h+2s \qquad(\nu(c)\ge-s,\ s\ge0 \text{ rational}).
\end{equation}
Indeed the generic value for \(c=t^{-s}u\) is convex in \(s\) by
one-inner-end comparisons. It starts at most \(h\) and is bounded
above by \(3h+2s\) throughout. Its secant slopes from zero are thus
at most \(2\). Specialization proves the assertion.

If \(d(B,L_z)=2h+2s\) with \(s\ge0\), apply
\eqref{eq:motion-midpoint} with \(c=(z-b)/v\). The triangle inequality
forces equality and \(m\in[B,L_z]\). In particular every leaf at
distance \(2h\) from \(B\) gives the same midpoint. For each rational
\(\ell\le1\), denote by \(P(z,\ell)\) the point of \([L_z,L_{z'}]\)
at distance \(1-\ell\) from \(L_z\) for any sufficiently distant
\(L_{z'}\). For \(\ell<1\) choose first \(z'\) at distance
\(2(1-\ell)\) and use the midpoint just proved. For \(\ell=1\)
take \(L_z\). Two distinct leaves \(L_z,L_{z'}\) at distance \(2h\)
meet at depth \(1-h=\nu(z-z')\). Choose a third leaf sufficiently
far from both. The midpoint test applied from each original leaf
places their common midpoint on both geodesics to the third. Above
the meeting depth the two paths lie in opposite halves of
\([L_z,L_{z'}]\). Below it they lie on the same segment from the
midpoint to the third leaf, by uniqueness of geodesics. Coincident
leaves have the same path at every depth. This gives
\begin{equation}\label{eq:motion-tree}
 d(P(z,\ell),P(z',\ell'))=
       \ell+\ell'-2\min\{\ell,\ell',\nu(z-z')\}.
\end{equation}
The point \(P(z,\ell)\) thus depends only on the disk
\(\{z':\nu(z'-z)\ge\ell\}\). Its ancestors correspond to the larger
disks containing it.

In a finite star comparison we can now prescribe points
\(P(a+bc,\ell)\) with depths harmonic across that star. Choose a
far leaf \(L_{t^{-N}}\) with \(N\) sufficiently large. Its distance
to every \(L_{a+bc}\) in the finite test is the same, and the
corresponding segment contains the prescribed site. The prescribed distances to these
two endpoints are \(1-\ell\) and its complement in their common
distance. Lemma~\ref{lem:geodesic} applies. It also gives specialization
for constant depths. In particular the price of a fixed nonzero tensor along a ray
\(c=\xi+t^\rho u\) is convex in \(\rho\) if each depth is affine in
\(\rho\), since any three rational radii are compared by a star with
one inner end at zero. At \(\rho=0\), the generic values obtained
from \(u\) and \(\xi+u\) coincide for integral \(\xi\), by the two
specializations in opposite directions.

\subsection{The price of a common label at four sites}

Place the common label \(\Delta_{[4]}(q)=\sum_{j=1}^q e_j^{\otimes4}\),
\(q\ge2\), at the tips of four equal branches of the disk tree. Choose
the first three centers \(z_1=0,z_2=1,z_3\in F\setminus\{0,1\}\)
and let the fourth center \(z_4=c\) vary. At a fresh \(c=u\), the
four tips \(P(z_i,\ell)\) at rational depth \(0<\ell<1\) have common
origin \(P(0,0)\) and pairwise distance \(2\ell\). Denote the price
of their label divided by \(\log q\) by \(D_\ell\). It vanishes at
\(\ell=0\). Joint geodesic convexity and the upper bound \(4\ell\)
from four size-\(q\) wires to the common origin give a finite
one-sided slope
\[
 \underline D=\lim_{\ell\downarrow0}\frac{D_\ell}{2\ell}.
\]
The denominator is the rate of a single wire between two of the tips.
We estimate the ratio by moving the fourth center toward each of the
three fixed centers in turn.

On the arm toward \(z_i\), let \(c=z_i+t^r u\) for \(0\le r\le2\ell\).
Increase the depths of \(i,4\) by \(r/2\) and decrease those of the
other pair by \(r/2\). Take \(2\ell\le1\), so all four depths remain
in the tree. By \eqref{eq:motion-tree} the two within-pair distances
are \(h=2\ell-r\), and all four distances between pairs are \(2\ell\).
The path joining the pair midpoints grows by precisely the amount
the within-pair distances shrink. Figure~\ref{fig:collision} depicts
the collision with \(z_1\). Relabeling gives the collision with \(z_3\)
used in Section~\ref{sec:boundary}. When \(r=2\ell\), both pairs have
merged. The grouped vectors \(e_j\otimes e_j\) are linearly independent
at each end, so the label is equivalent under local maps to one
size-\(q\) wire. Its rate at the collision is exactly \(2\ell\).

\begin{figure}[htbp]
\centering
\begin{tikzpicture}[
    x=1cm,y=1cm,
    outer/.style={draw=ink,line width=.75pt},
    trunk/.style={draw=sea,line width=1.15pt},
    site/.style={circle,fill=ink,inner sep=0pt,minimum size=4.2pt},
    branch/.style={circle,draw=ink,fill=white,inner sep=0pt,minimum size=4.0pt},
    lab/.style={font=\small,inner sep=2pt},
    stage/.style={font=\small,anchor=north,inner sep=1pt}
  ]
  \begin{scope}
    \coordinate (p) at (0,0);
    \coordinate (one) at (-.86,.66);
    \coordinate (four) at (-.86,-.66);
    \coordinate (two) at (.86,.66);
    \coordinate (three) at (.86,-.66);
    \draw[outer] (one)--(p)--(two) (four)--(p)--(three);
    \node[lab,above right] at (-.47,.36) {$\ell$};
    \node[branch] at (p) {};
    \foreach \pt/\val/\side in {one/1/left,four/4/left,two/2/right,three/3/right}{
      \node[site] at (\pt) {};
      \node[lab,\side=2pt] at (\pt) {$\val$};
    }
    \node[stage] at (0,-1.02) {$h=2\ell$};
  \end{scope}
  \draw[-{Stealth[length=4pt]},draw=black!50] (1.51,0)--(2.13,0);
  \begin{scope}[shift={(4.11,0)}]
    \coordinate (leftbranch) at (-.67,0);
    \coordinate (rightbranch) at (.67,0);
    \coordinate (one) at (-1.36,.66);
    \coordinate (four) at (-1.36,-.66);
    \coordinate (two) at (1.36,.66);
    \coordinate (three) at (1.36,-.66);
    \draw[outer] (one)--(leftbranch)--(four) (two)--(rightbranch)--(three);
    \draw[trunk] (leftbranch)--(rightbranch);
    \node[lab,above] at (0,.02) {$2\ell-h$};
    \node[lab,above right] at (-1.11,.41) {$h/2$};
    \node[branch] at (leftbranch) {};
    \node[branch] at (rightbranch) {};
    \foreach \pt/\val/\side in {one/1/left,four/4/left,two/2/right,three/3/right}{
      \node[site] at (\pt) {};
      \node[lab,\side=2pt] at (\pt) {$\val$};
    }
    \node[stage] at (0,-1.02) {$0<h<2\ell$};
  \end{scope}
  \draw[-{Stealth[length=4pt]},draw=black!50] (6.10,0)--(6.72,0);
  \begin{scope}[shift={(8.36,0)}]
    \draw[trunk] (-.90,0)--(.90,0) node[lab,midway,above] {$2\ell$};
    \node[site] at (-.90,0) {};
    \node[site] at (.90,0) {};
    \node[lab,below] at (-.90,-.05) {$1,4$};
    \node[lab,below] at (.90,-.05) {$2,3$};
    \node[stage] at (0,-1.02) {$h=0$};
  \end{scope}
\end{tikzpicture}
\caption{Collision of the fourth tip with the first. On the middle
tree the distance inside each pair is \(h\), and the segment between
pair midpoints has length \(2\ell-h\). At the right endpoint the
common label is a single wire of length \(2\ell\).}
\label{fig:collision}
\end{figure}


Across these three arms the depth slopes sum to \(-1/2\) at each of
the first three sites and to \(3/2\) at site \(4\). To make all depths
harmonic, add two arms of length \(\ell\). On an unused finite arm
\(a\in F\setminus\{z_1,z_2,z_3\}\) increase every depth at slope
\(1/2\). By specialization its end rate is at most \(D_{3\ell/2}\).
On the outer arm use \(c=t^{-r}u\). Only the fourth depth changes, at
slope \(-2\). At \(r=\ell\) this site is \(P(0,-\ell)\). Write
\(D_\ell^-\) for the resulting label rate.
The ray price functions are continuous at their origins. Indeed for
small signed \(r\) the same affine depth prescriptions and \(t^r\) or
\(t^{-r}\) give a convex extension across zero.

Figure~\ref{fig:harmonic} records the five depth changes in a single
comparison. The sum down each column is zero, which is the harmonicity
condition for the corresponding site.

\begin{figure}[htb]
\centering
\begin{tikzpicture}[font=\small,
 inner/.style={draw=ink,line width=.8pt},
 end/.style={inner sep=2pt,anchor=west},
 twig/.style={draw=sea,densely dashed,line width=.85pt},
 up/.style={draw=rust,line width=1.05pt}]
 \draw[inner] (0,0)--(1.9,1.4) node[end] at (1.96,1.4){$z_1$};
 \draw[inner] (0,0)--(2.32,.4) node[end] at (2.38,.4){$z_2$};
 \draw[inner] (0,0)--(2.18,-.9) node[end] at (2.24,-.9){$z_3$};
 \draw[twig] (0,0)--(.48,-1.08);
 \draw[up,{Stealth[length=4pt]}-] (-1.18,0)--(0,0);
 \node at (-1.58,0) {$\infty$};
 \node at (.58,-1.38) {$a$};
 \fill[ink] (0,0) circle (1.8pt);
 \node at (-.19,.36) {$u$};
 \node[anchor=west] at (3.58,0) {$
  \renewcommand{\arraystretch}{1.22}
  \begin{array}{c | c c c c | l}
     &1&2&3&4&\text{end estimate}\\ \hline

    z_1 &1&-1&-1&1& 2\ell\\
    z_2 &-1&1&-1&1& 2\ell\\
    z_3 &-1&-1&1&1& 2\ell\\
    a&1&1&1&1& D_{3\ell/2}\\
    \infty&0&0&0&-4&D_\ell^-
  \end{array}$};
\end{tikzpicture}
\caption{The five arms based at the generic point \(u\) of the
parameter disk. The three collision arms have length \(2\ell\), the
unused finite arm \(a\) and the outer arm have length \(\ell\). The
table gives twice the four depth slopes and an upper bound for the
label rate at each end.}
\label{fig:harmonic}
\end{figure}


We compare \(D_\ell^-\) to the generic rate. Start the outer ray
at a fixed small depth \(\lambda>0\) and stop at \(0<R<\lambda\),
with all lengths rational. Write \(D(\lambda,R)\) for the rate.
The first three sites have radius \(\lambda\) about \(P(0,0)\).
By \eqref{eq:motion-tree} the fourth site also has radius \(\lambda\),
with its first \(R\) units of geodesic going toward negative depths.
Contract all four radii from \(\lambda\) to \(\ell\le R\).
This gives the placement defining \(D_\ell^-\), so convexity gives
\[
   \frac{D_\ell^-}{2\ell}\le\frac{D(\lambda,R)}{2\lambda}.
\]
Taking the upper limit in \(\ell\), then \(R\downarrow0\) by ray
continuity, then \(\lambda\downarrow0\), yields
\(\limsup D_\ell^-/(2\ell)\le \underline D\). Substituting the upper bound for
the fourth finite end into \eqref{eq:motion-star} gives
\begin{equation}\label{eq:motion-fourstar}
 0\le 3\frac{2\ell-D_\ell}{2\ell}
      +\frac{D_\ell^--D_\ell}{\ell}
      +\frac{D_{3\ell/2}-D_\ell}{\ell}.
\end{equation}
The three collision arms contribute \(3(1-\underline D)\) in the limit. The outer
arm has upper limit at most zero, and the last term tends to \(\underline D\) since
\(D_{3\ell/2}/(3\ell)\to \underline D\). Hence \(0\le 3(1-\underline D)+\underline D\), giving
\begin{equation}\label{eq:motion-upper}
                        \underline D\le 3/2 .
\end{equation}

The rectangular comparison of Section~\ref{sec:boundary} will use this
bound near the endpoint \(h=0\) in Figure~\ref{fig:collision}. There
the three fixed tips form a tripod with a pair of short arms and one
long arm. Propagating the resulting limiting price equality to the
balanced tripod gives a common-label rate of \(8/3\). The three
matrix multiplication wires on that tripod have total length \(6\),
so the rank comparison has exponent \(6/(8/3)=9/4\).
